\section{Related Work}
\label{sec:related}

\vspace{4pt}
\noindent\textbf{LLMs for parallel and HPC code.}
A growing line of work studies whether language models can write parallel code.
\textsc{ParEval}~\cite{pareval} is the standard benchmark, evaluating
pass@$k$ correctness and speedup for model-generated OpenMP, MPI, CUDA, and
Kokkos, and documents a large gap between models' serial and parallel
competence; \textsc{HPC-Coder}~\cite{hpccoder} fine-tunes models on HPC code, and
empirical studies of model-generated OpenMP find frequent semantic
errors---missing reductions, wrong data scoping, absent
synchronization~\cite{openmpunderstand}. These efforts adjudicate correctness on a
\emph{single} execution at a fixed degree of parallelism. We build directly on
\textsc{ParEval}'s tasks and serial references, but change the object of study:
rather than scoring one run, we ask whether generated code is correct across the
\emph{configuration} axis---thread count and repetition---and whether an agent can
be made to repair the failures that this axis exposes. Distributed
data-processing pipelines of this kind are the pre- and post-processing stages of
AI-coupled HPC workflows~\cite{aicoupled}, where a silent correctness bug
propagates into the science.

\vspace{4pt}
\noindent\textbf{Verifiable-reward RL and performance-aware code generation.}
A growing body of work post-trains code models with reinforcement learning from
\emph{verifiable} rewards---an execution checker, not human preference, supplies the
signal (RLVR~\cite{rlvr}, optimized with GRPO~\cite{grpo}). Closest to us, several efforts
reward \emph{performance} directly for HPC code: \textsc{RLPF} aligns models to generate
faster code from measured speedup~\cite{rlpf}, online RL feeds real-machine benchmark
throughput back as the reward~\cite{rlhpc_realmachine}, and CUDA-kernel work combines
correctness, performance, and structural rewards~\cite{cudaperf}; \textsc{PCEBench}
benchmarks parallel code on both correctness and performance~\cite{pcebench}. We differ in
\emph{claim}, not in noticing that speed matters: we do \emph{not} propose optimizing a
performance reward. We show that the \emph{default} single-axis correctness reward---the
one an execute-and-fix agent already uses---is \emph{saturated} at the frontier and
\emph{gameable by construction} (a serial program is its optimum), and we \emph{measure}
the speedup a correctness-only objective discards. The two-axis reward and its guarded
performance gate are the remedy that follows from that diagnosis, not the contribution we
assert over prior performance-rewarded generation.

\vspace{4pt}
\noindent\textbf{Functional-correctness evaluation.}
Code-generation evaluation has centered on functional correctness of small,
self-contained problems---HumanEval~\cite{codex}, its rigorous
re-evaluation~\cite{humaneval_reliability}, and MBPP-style suites---and,
at repository scale, SWE-bench~\cite{swebench}. All adjudicate a single execution
against a fixed test suite. EvalPlus~\cite{humaneval_reliability} shows that code
passing a visible suite often fails under additional tests; our work is the
parallel analog along an execution axis developers do not test: the same program,
correct at one thread count, fails at another.

\vspace{4pt}
\noindent\textbf{Dynamic race detection.}
A mature line of HPC tools detects data races dynamically:
\textsc{ThreadSanitizer}~\cite{tsan} instruments memory accesses at scale, and
\textsc{Archer}~\cite{archer}, \textsc{ROMP}~\cite{romp}, and
\textsc{Sword}~\cite{sword} specialize happens-before race detection for
\textsc{OpenMP}. These are far more powerful race \emph{detectors} than our
output-differential sweep, which only observes end-to-end disagreement. We use them
as the reference point for what our estimator can and cannot see
(Section~\ref{sec:secondaxis}): our contribution is not a better race detector but
the empirical finding that, for frontier-model output on idiomatic tasks, the race
they hunt is largely absent, and the binding constraint has moved to whether the
accepted code is parallel at all. Our robustness estimator is also the
parallel-code instance of the general \emph{flaky-test} problem---tests whose
verdict depends on nondeterministic execution~\cite{flakytests}---repetition being
our defence against it.

\vspace{4pt}
\noindent\textbf{Differential and metamorphic testing.}
Executing a program under multiple configurations and flagging disagreements is
the essence of differential testing~\cite{mckeeman}; checking properties
preserved under transformations is metamorphic testing~\cite{metamorphic}. In
data systems this lineage is rich: query synthesis and reference-engine
construction have found many engine bugs~\cite{rigger2020pqs,rigger2020norec}, and
\textsc{SparkFuzz}~\cite{sparkfuzz} differentially tests a query engine against a
single-node reference. Our verifier applies the differential lens along the
\emph{parallelism} axis, and---unlike these offline fuzzers---packages it as a
tool an \emph{agent} invokes inside its generation loop: the object under test is
the model's program, not the engine.

\vspace{4pt}
\noindent\textbf{Self-repairing and agentic code generation.}
A large body of work lets models improve their output through execution feedback:
Self-Debugging feeds runtime traces back to the model~\cite{selfdebug}, Reflexion
adds verbal self-critique~\cite{reflexion}, and Self-Refine iterates on
self-generated feedback~\cite{selfrefine}; in HPC, LASSI builds a
compile-and-run self-correction loop for parallel-code
translation~\cite{lassi}. These loops are only as good as the signal they consume.
Our tool-ablation isolates exactly this signal quality: in principle an
execute-and-fix loop with a single-configuration signal cannot see the
nondeterministic parallel bugs that a \emph{differential} signal would catch, so we
hold the loop fixed and vary only the verifier---and report, as an honest empirical
finding, when that distinction does and does not matter. To our knowledge no prior
work packages a distribution-aware differential verifier as an agent tool for
parallel-code correctness, nor measures the parallel-correctness gap and its
agentic repair on a standard HPC benchmark.
